%--------------------------------------
% Caso de Uso: Pagar Reservación
%--------------------------------------

\begin{UseCase}{CU-20}{Pagar Reservación}{
	Permitir que un \hyperlink{Propietario}{Propietario} o \hyperlink{Administrador}{Administrador} registre un pago asociado a una reservación.
}
	\UCitem{Versión}{\color{Gray}1.0}
	\UCitem{Autor}{\color{Gray}Equipo de Desarrollo}
	\UCitem{Supervisa}{\color{Gray}Coordinador del Sistema}
	\UCitem{Actor}{\hyperlink{Propietario}{Propietario} o Recepcionista}
	\UCitem{Propósito}{Registrar pagos para cubrir total o parcialmente el monto de una reservación.}
	\UCitem{Entradas}{Monto a pagar, método de pago y referencia opcional.}
	\UCitem{Origen}{Pantalla de detalle de reservación}
	\UCitem{Salidas}{Pago registrado, saldo actualizado y comprobante generado.}
	\UCitem{Destino}{Pantalla de detalle de reservación y base de datos}
	\UCitem{Precondiciones}{
		\begin{itemize}
			\item La reservación debe existir.
			\item La reservación debe tener saldo pendiente.
			\item El monto ingresado debe ser válido.
			\item El método de pago debe estar disponible.
		\end{itemize}
	}
	\UCitem{Postcondiciones}{
		\begin{itemize}
			\item El pago queda registrado en el sistema.
			\item El saldo de la reservación se actualiza.
			\item Se genera un comprobante de pago.
			\item El estado de la reservación puede actualizarse según el monto pagado.
		\end{itemize}
	}
	\UCitem{Errores}{
		Monto inválido, método de pago no disponible, error en la transacción, pago duplicado.
	}
	\UCitem{Tipo}{Caso de uso primario}
	\UCitem{Observaciones}{
		\begin{itemize}
			\item Se permiten pagos parciales o totales.
			\item Existen distintos métodos de pago.
			\item Los pagos en línea requieren validación externa.
		\end{itemize}
	}
\end{UseCase}

%--------------------------------------
\begin{UCtrayectoria}
	\UCpaso[\UCactor] Accede a la pantalla de detalle de una reservación.
	\UCpaso[\UCactor] Selecciona la opción para registrar un pago.
	\UCpaso El sistema muestra la información de la reservación y el saldo pendiente.
	\UCpaso El sistema muestra los métodos de pago disponibles.
	\UCpaso[\UCactor] Ingresa el monto a pagar.
	\UCpaso El sistema valida que el monto sea correcto.
	\UCpaso[\UCactor] Selecciona el método de pago.
	\UCpaso El sistema procesa el pago según el método seleccionado.
	\UCpaso[\UCactor] Confirma la operación.
	\UCpaso El sistema registra el pago.
	\UCpaso El sistema actualiza el saldo de la reservación.
	\UCpaso El sistema evalúa si el pago cubre el total de la reservación.
	\UCpaso El sistema genera el comprobante de pago.
	\UCpaso El sistema muestra un mensaje indicando que el pago fue registrado correctamente.
	\UCpaso[\UCactor] Puede imprimir o guardar el comprobante.
	\UCpaso El sistema actualiza la información mostrada de la reservación.
\end{UCtrayectoria}

%--------------------------------------
\begin{UCtrayectoriaA}{A}{Monto inválido}
	\UCpaso El sistema informa que el monto ingresado no es válido.
	\UCpaso El flujo regresa a la captura del monto.
\end{UCtrayectoriaA}

%--------------------------------------
\begin{UCtrayectoriaA}{B}{Pago rechazado}
	\UCpaso El sistema informa que la transacción no pudo completarse.
	\UCpaso El flujo regresa a la selección del método de pago.
\end{UCtrayectoriaA}

%--------------------------------------
\subsection{Puntos de extensión}
\UCExtenssionPoint{
	El usuario desea configurar un esquema de pagos a plazos.
}{
	Del paso 5.
}{
	Configurar Plan de Pagos.
}

%--------------------------------------
% Fin del Caso de Uso
